\section{Open Questions (5 concrete follow-ups)}

The following questions are genuinely open at the end of this replication --- each has
a concrete numerical or analytical next step that a follow-on wave can execute without
paid endpoints.

\begin{enumerate}
  \item \textbf{Higher-order Suzuki formulas.}
  Do the paper's predicted slopes hold at $p=4$ (Yoshida) and $p=6$ (Suzuki) on the
  same TFIM and Hubbard-dimer instances, and does the state-dependent bound continue
  to be tight to the same $\sim 1$--$2\%$ level as the order grows?
  \emph{Next step:} extend \texttt{code/trotter\_strang.py} and \texttt{code/hubbard\_dimer.py}
  to compose the Yoshida-4 and Suzuki-6 product formulas (well-known coefficient sets),
  sweep $r \in \{2,\ldots,64\}$ at $t=1.0$, fit log-log slopes, compare to $-4$ and $-6$
  in both operator-2-norm and state error. Runtime $<10$ s CPU.

  \item \textbf{Tightness of the state-dependent prefactor vs Childs--Su--Tran--Wiebe.}
  Can the paper's state-dependent (strong) bound be quantitatively separated from the
  standard Childs--Su--Tran--Wiebe commutator bound on a chosen state that makes the
  state-dependent constant small?
  \emph{Next step:} for the same TFIM instance, compute (a) the standard sum-of-nested-commutator-norms
  prefactor and (b) the paper's state-dependent prefactor evaluated on $|+\rangle^{\otimes n}$
  and on a low-energy eigenstate; report the ratio and confirm it is $\ge 1$, strictly $>1$
  on the low-energy state.

  \item \textbf{Integration with LCU / qubitization.}
  How does the state-dependent bound behave when integrated with LCU or qubitization
  primitives, where the cost model is queries to a block-encoding of $H$ rather than the
  number of Trotter steps $r$?
  \emph{Next step:} take a small qubitized-simulation benchmark (e.g.\ the 4-site Fermi-Hubbard
  block-encoding of Babbush et al.\ 2018), compare the raw block-encoding cost against a
  Trotter-refined estimate that plugs in the paper's state-dependent constant, at fixed
  $t$ and target error $10^{-3}$.

  \item \textbf{Chemistry-scale instance.}
  Does the paper's slope prediction remain within the same $\sim 1$--$2\%$ band on a
  genuine second-quantized chemistry Hamiltonian (H$_4$ chain STO-3G, LiH minimal basis,
  or 2$\times$4 Fermi-Hubbard) rather than the 4-orbital Hubbard dimer used here?
  \emph{Next step:} use PySCF/OpenFermion to build the H$_4$-chain second-quantized
  Hamiltonian in STO-3G, apply Jordan--Wigner or Bravyi--Kitaev, split into hopping ($A$)
  and Coulomb ($B$), re-run the $r$-sweep and slope-fit protocol at $t=1.0$.

  \item \textbf{Sample-cost Pareto frontier.}
  What is the Pareto frontier between (a) increasing Trotter order $p$, (b) increasing
  step count $r$, and (c) improving the state-dependent constant via better state preparation,
  at fixed target infidelity for the paper's target regime?
  \emph{Next step:} for the TFIM instance at $t=1.0$ and target error $10^{-3}$,
  plot gate count required to reach the target as a function of $(p,r)$ for
  $p \in \{1,2,4,6\}$ with state fixed to $|+\rangle^{\otimes n}$; then repeat with a
  variationally-prepared low-energy state and see whether the state-dependent improvement
  shifts the Pareto frontier meaningfully.
\end{enumerate}
